% Promoted after complete independent root and reduction/descent proof audits. No Petri/variation nonvanishing assertion.
\subsection{A norm pencil on the determinant curve}

\begin{lemma}[The determinant cokernel is residual to a squared pencil]
\label{lem:DE-A3-positive-pencil-candidate}
\status{Proved; complete independent root and reduction/descent proof audits passed.}
Keep the algebraically closed characteristic-two hypotheses and notation of
Lemma~\ref{lem:DE-A3-positive-determinant-curve}, and take a general pair
with disjoint zero schemes. On $D=D_s$ the images of $s_0$ and $s_1$
generate a line subbundle $R_0\subset E_0|_D$, after viewing the latter
section as a map $P^{-1}|_D\to E_0|_D$. Then
\[
 R_0\simeq\mathcal O_D(Z_0),\qquad
 R_0P\simeq\mathcal O_D(Z_1),\qquad
 P|_D\simeq\mathcal O_D(Z_1-Z_0).
\]
In particular $Z_0,Z_1$ are effective Cartier divisors on $D$, including
when $D$ is nonreduced. Put $A=R_0^{\otimes2}$. The squares of the two
sections, using the fixed trivialization of $P^2$, give a base-point-free
pencil $\langle a_0,a_1\rangle\subset H^0(D,A)$ with divisors $2Z_0,2Z_1$.
The line bundle in (DE-TB1) is
\[
                         R_s\simeq\omega_D\otimes A^{-1}.
                                                        \tag{DE-TP1}
\]
Here $\deg A=b$ and $p_a(D)=b+1$.
\end{lemma}
\begin{proof}
At every point of $D$, at least one of $s_0,s_1$ is nonzero, and their
determinant vanishes. In a frame with the nonvanishing section first,
their matrix has a unit entry and determinant zero. Its image on $D$
is therefore a line subbundle; this constructs $R_0$ even over the
nonreduced local rings of $D$. Near $Z_0$, use $s_1$ to trivialize
the common line. The two equations of $Z_0$ on the smooth surface can
then be written $(f,w)$, where $w$ cuts out $D$. They form a regular
sequence. Consequently $f$ is a nonzerodivisor on $D$ and cuts out
$Z_0$ there. The same argument treats $Z_1$. The asserted line-bundle
identifications follow from their sections. Their squares have no common
zero, since $Z_0\cap Z_1=\varnothing$.

Tensor (DE-TB1) by $\mathcal O_D$. Since $D$ is Cartier, its initial
Tor term is $R_s\otimes\mathcal O_D(-D)$, and hence
\[
 R_s(-D)=\ker\bigl(\mathcal O_D^2\longrightarrow Q_s|_D\bigr).
\]
The earlier local chart has columns $(1,0)$ and $(u f^2,w)$,
where $e^2\mapsto u$ for the chosen frame $e$ of $P$.
On $D$ their kernel is the kernel of the pencil row $(1,u f^2)$.
Globally this row is $(a_0,a_1):\mathcal O_D^2\to A$, whose kernel
is $A^{-1}$. Thus $R_s(-D)=A^{-1}$. Adjunction on the abelian surface
identifies $\mathcal O_D(D)$ with $\omega_D$, proving (DE-TP1).
The degree statements follow from $\operatorname{length}Z_i=b/2$ and
the previously computed arithmetic genus. No reducedness or smoothness
of $D$ was used.
\end{proof}

\begin{lemma}[Norm factorization and the actual ordinary differential]
\label{lem:DE-A3-positive-norm-pencil-candidate}
\status{Proved; complete independent root and reduction/descent proof audits passed.}
With the same general pair, the norm of the $\alpha_2$-quotient
$\nu:C\to T$ gives a section
\[
                   \sigma\in H^0(D,R_sP)
\]
which is a nonzerodivisor on $D$. If $S_i\in H^0(T,L)$ is the
Frobenius-twisted section corresponding to $s_i$, then
\[
                         S_i|_D=\sigma a_i.             \tag{DE-TP2}
\]
There is a canonical isomorphism
\[
 \frac{H^0(T,L)}{\langle S_0,S_1\rangle}
 \simeq
 \frac{H^0(D,\omega_DP)}{\sigma\langle a_0,a_1\rangle}.
                                                        \tag{DE-TP3}
\]
The fixed trivialization of $P^2$ also defines a nonzero invariant
differential $\eta_P\in H^0(T,\Omega_T^1)$ with
$C(\eta_P)=\eta_P$. Wherever $a_0a_1$ is invertible on $D$,
\[
                      d\log(a_1/a_0)=\eta_P|_D.         \tag{DE-TP4}
\]
This identity does not assert that its right side is nonzero on $D$.
\end{lemma}
\begin{proof}
Locally the finite flat $\alpha_2$-algebra is
$\mathcal O_T[z]/(z^2-h)$. In a trivialization of $J$, its norm is
\[
                         n(a+bz)=a^2+h b^2.
\]
Determinants of the transition matrices identify its target with
$\det\nu_*J=L$; the determinant of $\nu_*\mathcal O_C$ is trivial,
as the torsor algebra is an extension of $\mathcal O_T$ by itself.
Thus this construction gives the intrinsic quadratic norm
$n:E_0\to L$. For $E_1=E_0P$, the same norm is multiplied by the
local unit $u$ specifying $P^2=\mathcal O_T$. This is precisely the
descent of the square of the original section through relative
Frobenius, so $n(s_i)=S_i$.

The restriction of the norm to the line $R_0$ is a homomorphism
$A=R_0^2\to L|_D$. It is therefore a section of
$L|_D\otimes A^{-1}=R_sP$, using $\omega_D=(LP)|_D$.
This is $\sigma$, and homogeneity gives (DE-TP2). For a general
pair, the divisors of the sections of $J$ and $JQ$ on $C$ have no
common component: both line bundles are globally generated, and
one may choose the second section to avoid the finitely many
components of the first. Their norm divisors on $T$ likewise have
no common component, since $\nu$ is a homeomorphism on underlying
spaces. If $\sigma$ belonged to a minimal prime of $D$, (DE-TP2)
would put both norm sections in the corresponding prime. This is
impossible. The Cartier divisor $D$ is Cohen--Macaulay and has no
embedded associated primes, so $\sigma$ is a nonzerodivisor.

The restriction sequence is
\[
                 0\longrightarrow P\longrightarrow L
                    \longrightarrow L|_D\longrightarrow0.
\]
Acyclicity of $P$ identifies $H^0(T,L)$ with $H^0(D,L|_D)$.
As $L|_D=\omega_DP$, equation (DE-TP2) gives (DE-TP3).

Finally, write $e^2\mapsto u$ in any local frame of $P$.
On overlaps $u$ changes by the square of a transition unit,
so $d\log u$ glues to a global differential $\eta_P$.
The usual Cartier identity gives $C(d\log u)=d\log u$.
If this differential were zero, the Cartier kernel calculation
on the smooth perfect-base surface would make every $u$ locally
a square. The uniquely determined square roots then normalize
all transition functions to one, trivializing $P$, a contradiction.
Every global differential on an abelian variety is invariant.
In the determinant chart the pencil ratio is exactly $u f^2$,
so its logarithmic differential is $d\log u$. This proves
(DE-TP4), with no assumption that $u$ itself has a local square root.
\end{proof}

\begin{lemma}[The entire Cartier contraction is a curve extension]
\label{lem:DE-A3-positive-pencil-contraction-candidate}
\status{Proved; complete independent root and reduction/descent proof audits passed.}
The two bundles are orthogonal also in cohomology:
\[
                  \operatorname{Ext}^q_T(E_0,E_1)=0
                         \quad\text{for every }q.
                                                        \tag{DE-TP5}
\]
There is an exact sequence of bundles on $D$
\[
 0\longrightarrow A\longrightarrow G_s:=E_0|_D\otimes R_0
                  \longrightarrow L|_D\longrightarrow0,
                                                        \tag{DE-TP6}
\]
whose global-section map, under the natural identifications
\[
 H^0(D,G_s)=H^0(T,E_0\oplus E_1),\qquad
 H^0(D,L|_D)=H^0(T,L),
\]
is exactly the Cartier contraction $C_s$. Moreover $H^1(D,G_s)=0$.
Thus, for a sharp source,
\[
                  h^0(D,A)=h^1(D,A)=h^0(D,R_s)=2,
\]
and the two pencil sections form a basis of $H^0(D,A)$.

Let $\beta\in H^1(T,\mathcal O_T)$ be the class of the
$\alpha_2$-torsor $\nu$. If $\partial$ is the boundary of
(DE-TP6), then
\[
                    \partial(\sigma a)=\beta|_D\cup a
                         \quad(a\in H^0(D,A)).          \tag{DE-TP7}
\]
For a sharp source this cup-product map $H^0(D,A)\to H^1(D,A)$
is an isomorphism. The scalar Frobenius twists in identifying
the original kernel sections with their squared pencil sections
do not affect this isomorphism assertion.
\end{lemma}
\begin{proof}
The self-fibre product of the $\alpha_2$-torsor has square-zero
diagonal ideal isomorphic to $\mathcal O_C$. Consequently
$\nu^*\nu_*J$ has a two-step filtration with both quotients $J$.
Projection formula and finite flat base change identify the cohomology
computing (DE-TP5) with that of
\[
                  (\nu^*\nu_*J)^\vee\otimes JQ.
\]
This bundle has a filtration with both quotients $Q$. Since the
nontrivial degree-zero line bundle $Q$ is acyclic on $C$, the
whole cohomology vanishes. This proves (DE-TP5).

The two original sections give an exact sequence on the surface
\[
 0\longrightarrow\mathcal O_T\oplus P\longrightarrow E_0
       \longrightarrow i_*\bigl(L|_D\otimes R_0^{-1}\bigr)
       \longrightarrow0.
\]
For the cokernel identification, its Tor kernel on $D$ is
$P R_0^{-1}$, and tensoring back by $\mathcal O_D(D)=\omega_D$
gives $L|_D R_0^{-1}$. Dualizing the surface sequence gives
\[
 0\longrightarrow E_0^\vee\longrightarrow\mathcal O_T\oplus P
        \longrightarrow i_*(R_0P)\longrightarrow0.
\]
Here the last map is the row $(s_1,s_0)$, with the fixed
$P^2$ identification understood. Tensor this sequence by $E_1$.
Its left term is $E_0^\vee E_1$, acyclic by (DE-TP5), and the
other terms are $E_1\oplus E_0$ and $i_*G_s$. This gives the
claimed global-section identification and $H^1(G_s)=0$.
The inclusion of $R_0$ in $E_0|_D$, tensored by $R_0$, gives
(DE-TP6). Under the preceding map its quotient row is
$\beta(s_0,-)+\beta(s_1,-)$: this can be read in the determinant
chart, or directly by taking the wedge with the common line.
It is exactly $C_s$.

The long cohomology sequence therefore identifies $H^0(A)$
with $\ker C_s$ and $H^1(A)$ with $\operatorname{coker}C_s$.
At a sharp source these have dimension two. The two known
kernel vectors are $s_0,s_1$, and under the same identifications
they are $a_0,a_1$. Serre duality on the Gorenstein Cartier curve,
together with (DE-TP1), gives $h^0(R_s)=h^1(A)=2$.

For the last assertion choose local torsor coordinates $z_i$
with $z_j-z_i=c_{ij}$, representing $\beta$. Multiplication by
$z_i$ is a local endomorphism of $E_0=\nu_*J$; this statement
is independent of the local trivialization of $J$.
If $r$ trivializes $R_0$ on $D$, then
\[
                   r\wedge(z_i r)=n(r)=\sigma r^2.
\]
Thus, for $a=f r^2$, the section $f(z_i r)\otimes r$ of $G_s$
lifts $\sigma a$. Its transition difference is $c_{ij}a$,
which proves (DE-TP7). Finally, first rank equal to stable rank
$b-2$ is equivalent to the Frobenius-twisted kernel span
$\langle S_0,S_1\rangle$ mapping isomorphically to the cokernel
of $C_s$. By (DE-TP2) and (DE-TP7), this is precisely the
invertibility of cup product by $\beta$ on $H^0(A)$.
\end{proof}

\begin{corollary}[An exact formula for the full first rank of a pullback]
\label{cor:DE-A3-positive-first-rank-candidate}
\status{Proved; complete independent root and reduction/descent proof audits passed.}
Let $s$ be a sharp source, and let $\pi':\widetilde T\to T$ be the
etale double cover occurring in the earlier variation reduction, of
class $\alpha\in H^1(T,\mathcal O_T)$. Then the full coherent
Cartier contraction at the pulled-back section has kernel dimension
\[
        2+\dim\ker\bigl(\alpha|_D\cup-:
                     H^0(D,A)\longrightarrow H^1(D,A)\bigr).
                                                        \tag{DE-TP8}
\]
Its first rank is accordingly
$2b-2-\dim\ker(\alpha\cup-)$, which may be $2b-4$, $2b-3$,
or $2b-2$. This statement concerns the full contraction;
the previously proved first and stable ranks of its restriction
to the toral image lattice remain $2b-4$.
\end{corollary}
\begin{proof}
All determinant-curve constructions commute with etale base change.
The pullback of $E_0^\vee E_1$ remains acyclic: pushforward along
$\pi'$ is an extension of two copies of this acyclic bundle.
The same argument applies to $P$, and the higher cohomology
of the pulled-back $E_i$ vanishes. Thus the proof of (DE-TP6)
identifies the pulled-back contraction kernel with
$H^0(\widetilde D,\pi'^*A)$.
The torsor trace sequence, tensored by $A$ and restricted to $D$,
has connecting homomorphism $\alpha|_D\cup-$. Its first three
cohomology terms give
\[
 0\longrightarrow H^0(D,A)\longrightarrow
 H^0(\widetilde D,\pi'^*A)\longrightarrow
 \ker\bigl(\alpha\cup-:H^0(D,A)\to H^1(D,A)\bigr)
 \longrightarrow0.
\]
The preceding lemma gives $h^0(A)=2$, proving (DE-TP8).
The full source and target spaces each have dimension $2b$.
\end{proof}

\begin{remark}[What this model has not established]
\label{rem:DE-A3-positive-norm-pencil-scope-candidate}
\status{Proved scope statement; the quadratic variation remains unresolved.}
The two lemmas identify the quotient domain of the variation with an
explicit twisted canonical space on a Cartier curve and retain both
the norm factor and the ordinary logarithmic differential. They do
not identify the matrix-valued variation with a Petri or residue map.
The preceding two additional statements now prove $h^0(A)=2$
for a sharp source and identify the two torsor cup products exactly.
They still give no smoothness of $D$, no generic separability of its
pencil, and no lower bound of three on the variation-image dimension.
The retained finite probe in fact has total image dimension two:
the one-dimensional first-weight image is contained in the
two-dimensional second-weight image, and their displayed generators
have an invertible sum. The common-kernel/common-image exclusion
for arbitrary mixed surfaces remains the open step. The later
trace-zero theorem, Proposition~\ref{prop:DE-A3-positive-cup-trace-zero-candidate},
reduces a possible failure to one nilpotent matrix line; it does
not supply the remaining image-dimension bound.
\end{remark}
